% ====================================================================
% Capítulo 2: El debate experimental clásico
% ====================================================================
\chapter{El debate experimental clásico: del agua hirviendo al sitio de nucleación controlado}
\label{ch:debate_clasico}

\section{La experiencia primigenia de Mpemba y Osborne (1969)}

Mpemba y Osborne~\cite{mpemba1969} realizaron una secuencia de experimentos con dos 
muestras de líquido en vasos abiertos colocados simultáneamente en un congelador 
doméstico. Sus protocolos exploraron tres variables independientes:
\begin{enumerate}
\item Temperatura inicial $\Tinit$ de la muestra.
\item Pureza y composición (agua destilada, leche, agua del grifo).
\item Material y geometría del recipiente.
\end{enumerate}

\noindent Sus observaciones, replicadas en \SI{50}{} pares de experimentos, mostraron 
que el agua inicialmente a $\SI{80}{\celsius}$ comenzaba a formar cristales en 
$\SI{30}{\minute}$ aproximadamente, mientras que el agua inicialmente a 
$\SI{20}{\celsius}$ requería $\SI{100}{\minute}$. La métrica empleada fue el 
\emph{tiempo de comienzo de la cristalización}, definido por la formación de la 
primera capa visible de hielo en la superficie superior del líquido.

\section{Problemas epistemológicos de la definición clásica}

La definición operacional de Mpemba--Osborne sufre tres limitaciones que han 
contaminado más de cinco décadas de debate experimental~\cite{burridge2016}:

\begin{enumerate}
\item \textbf{Métrica subjetiva.} El ``comienzo de la cristalización'' es difícil 
de detectar de manera reproducible: depende de la iluminación, el tamaño de los 
cristales detectables y la presencia de núcleos heterogéneos en la pared del 
recipiente.
\item \textbf{Estocasticidad de la nucleación.} La cristalización del agua a 
$\SI{0}{\celsius}$ es un proceso estocástico que requiere superar una barrera 
energética de Helmholtz para formar un núcleo crítico. La distribución de 
tiempos de espera es exponencial, con varianza comparable a la media. Esto 
implica que las muestras pueden permanecer en estado sobreenfriado a 
$\SI{-10}{\celsius}$ durante decenas de minutos antes de cristalizar.
\item \textbf{Confusión entre enfriamiento y congelación.} El proceso global 
incluye dos etapas: (a) enfriamiento del líquido desde $\Tinit$ hasta el punto 
de fusión y eventual sobreenfriamiento, y (b) liberación del calor latente 
durante la cristalización. El efecto reportado puede deberse a cualquiera de 
las dos por separado.
\end{enumerate}

\section{El análisis crítico de Burridge y Linden (2016)}

Burridge y Linden~\cite{burridge2016} propusieron una redefinición rigurosa:

\begin{definition}[Efecto Mpemba estricto, Burridge--Linden 2016]
\label{def:mpemba_burridge}
Sean dos muestras de igual masa, igual recipiente y baño térmico común. Sean 
$E_{H}$ y $E_{C}$ las entalpías específicas requeridas para enfriar y congelar 
las muestras inicialmente caliente y fría desde sus condiciones iniciales hasta 
un estado final común. Sean $Q_{H} = E_{H}/t_{H}$ y $Q_{C} = E_{C}/t_{C}$ las 
tasas medias de transferencia de calor durante el proceso. El efecto Mpemba 
se considera observado si
\begin{equation}
    \frac{Q_{H}}{Q_{C}} > \frac{E_{H}}{E_{C}}.
    \label{eq:burridge_criterion}
\end{equation}
\end{definition}

\noindent Esta condición es equivalente a $t_{C} > t_{H}$, es decir, el tiempo de 
congelación completa de la muestra fría supera al de la caliente. La 
\cref{eq:burridge_criterion} se denomina la \emph{línea de Mpemba}.

Aplicando esta condición a los datos disponibles, Burridge y Linden concluyeron 
en 2016: \emph{``el efecto Mpemba no es observable en agua de manera 
estadísticamente significativa''.} Esta conclusión fue, sin duda, una llamada de 
atención para la comunidad: los reportes positivos del efecto se atribuyeron a 
sesgo de selección de datos, varianza de la nucleación o diferencias incontroladas 
en las condiciones iniciales.

\section{La rehabilitación experimental: Hallstadius y Burridge (2020)}

Tras el escepticismo de 2016, los mismos autores buscaron condiciones 
experimentales en las que el efecto pudiera reproducirse de manera 
sistemática~\cite{hallstadius2020}. Su hipótesis fue que el cuello de botella en 
la observación del efecto Mpemba no es la fase de enfriamiento del líquido, sino 
la fase de nucleación: una muestra que comienza a congelarse 
\emph{deterministamente} antes que otra disipará el calor latente antes y 
alcanzará el estado final de hielo primero.

El protocolo propuesto consistió en introducir asimetría \emph{sistemática} en la 
densidad de sitios de nucleación heterogéneos:

\begin{itemize}
\item Recipiente A (muestra inicialmente caliente): pared interior lijada con 
papel de esmeril P\SI{60}{} (rugosidad alta).
\item Recipiente B (muestra inicialmente fría): pared interior lisa, idéntica en 
geometría exterior, material y volumen.
\end{itemize}

\noindent Con este protocolo, en 12 pares de experimentos con muestras de 
\SI{200}{\milli\liter} de agua deionizada:

\begin{itemize}
\item La muestra inicialmente caliente cristalizó antes en \SI{11}{} de las 12 
parejas (probabilidad de hipótesis nula $p < 0.005$).
\item Se cumplió la condición de Burridge--Linden estricta 
(\cref{eq:burridge_criterion}) con $\Delta T_{0} = T_{H} - T_{C}$ tan grande como 
\SI{50}{\celsius}, correspondiente a una diferencia de entalpía 
$E_{H}/E_{C} \approx 1.5$.
\end{itemize}

\noindent La conclusión: el efecto Mpemba es real, pero \emph{condicionado} a la 
existencia de asimetrías en la cinética de nucleación.

\section{Mecanismos microscópicos propuestos en la literatura clásica}

Antes de la formulación markoviana de 2017, los mecanismos propuestos para el 
efecto Mpemba en agua incluían~\cite{burridge2016, hallstadius2020}:

\begin{enumerate}
\item \textbf{Evaporación}. El agua caliente pierde más masa por evaporación, 
reduciendo la masa que debe ser enfriada. Modelos cuantitativos predicen una 
contribución del 1--5\%, insuficiente para explicar las observaciones reportadas.
\item \textbf{Convección térmica}. Gradientes de densidad en el agua caliente 
generan convección que acelera el transporte de calor. Vynnycky y Maeno 
estimaron que esto contribuye en menos del 10\%~\cite{burridge2016}.
\item \textbf{Gases disueltos}. La solubilidad del aire en agua disminuye con la 
temperatura. El agua caliente contiene menos gas disuelto, lo que altera la 
cinética de cristalización heterogénea.
\item \textbf{Sobreenfriamiento diferencial}. El agua fría tiende a sobreenfriarse 
más profundamente antes de cristalizar, retrasando la liberación de calor 
latente.
\item \textbf{Memoria de enlaces de hidrógeno y supersolidez de piel}: el modelo 
de Zhang \emph{et al.}~\cite{zhang2014}, que se desarrolla en el 
\cref{ch:mecanismos_agua}.
\end{enumerate}

\section{Implicaciones para la teoría general}

El debate experimental clásico llevó a una bifurcación en la literatura:

\begin{enumerate}
\item Por un lado, los \emph{mecanicistas del agua}, que buscan en la estructura 
molecular de $\mathrm{H_2O}$ y sus enlaces de hidrógeno la explicación 
material del efecto.
\item Por otro lado, los \emph{formalistas markovianos}, quienes asumen que el 
efecto es genérico de los sistemas de no-equilibrio markovianos y prescinden de 
los detalles del medio.
\end{enumerate}

\noindent Ambos enfoques son válidos y complementarios. El presente documento los 
trata por separado en los \cref{ch:mecanismos_agua,ch:markoviano}, respectivamente, 
y los unifica en la perspectiva de termomayorización del 
\cref{ch:termomajorizacion}.

\begin{table}[h]
\centering
\caption{Resumen cronológico de los hitos experimentales clásicos en el debate del 
efecto Mpemba en agua.}
\label{tab:historial_experimental}
\small
\begin{tabularx}{\textwidth}{lll>{\raggedright\arraybackslash}X}
\toprule
\textbf{Año} & \textbf{Autores} & \textbf{Publicación} & \textbf{Contribución} \\
\midrule
350 a.C. & Aristóteles & Meteorología & Primera observación documentada \\
1620 & F. Bacon & Novum Organum & Confirmación cualitativa \\
1637 & R. Descartes & Les Météores & Confirmación cualitativa \\
1775 & J. Black & Edinburgh Trans. & Primer estudio sistemático \\
1969 & E. B. Mpemba, D. G. Osborne & Phys. Educ. \textbf{4}, 172 & Redescubrimiento moderno \\
1995 & D. Auerbach & Am. J. Phys. \textbf{63}, 882 & Rol del sobreenfriamiento \\
2009 & J. I. Katz & Am. J. Phys. \textbf{77}, 27 & Mecanismo de solutos \\
2016 & H. C. Burridge, P. F. Linden & Sci. Rep. \textbf{6}, 37665 & Análisis crítico, definición rigurosa \\
2020 & O. Hallstadius, H. C. Burridge & Proc. R. Soc. A \textbf{476}, 20190829 & Reproducibilidad por nucleación controlada \\
\bottomrule
\end{tabularx}
\end{table}
